\lecture{25}{Метод чередующихся направлений множителей}

\sourcevideo
    {https://www.youtube.com/playlist?list=PLOzRYVm0a65fhKDS8cYIC7YCuVGJ4FOJx}
    {Week 7: Lecture 25: ADMM Algorithm}

Метод чередующихся направлений множителей (ADMM) применяется к
задачам, в которых целевая функция разделяется по блокам, а переменные
связаны линейным равенством. Метод сохраняет квадратичный штраф
дополненной функции Лагранжа, но вместо совместной минимизации по всем
переменным выполняет последовательные блочные шаги.

\section{Двухблочная задача и дополненная функция Лагранжа}

Рассмотрим выпуклую задачу
\[
\begin{aligned}
    \min_{x\in\mathbb R^n,\,z\in\mathbb R^p}\quad
    &f(x)+g(z),
    \\
    \text{при условии}\quad
    &Ax+Bz=c,
\end{aligned}
\]
где
\[
    A\in\mathbb R^{m\times n},
    \qquad
    B\in\mathbb R^{m\times p},
    \qquad
    c\in\mathbb R^m.
\]
Прямая невязка равна
\[
    r(x,z)=Ax+Bz-c.
\]
Для двойственной переменной \(y\in\mathbb R^m\) обычная функция
Лагранжа имеет вид
\[
    \mathcal L(x,z,y)
    =
    f(x)+g(z)+y^{\mathsf T}(Ax+Bz-c).
\]
При параметре \(\rho>0\) дополненная функция Лагранжа определяется
формулой
\[
\begin{aligned}
    \mathcal L_\rho(x,z,y)
    ={}&
    f(x)+g(z)+y^{\mathsf T}(Ax+Bz-c)
    \\
    &+
    \frac{\rho}{2}\lVert Ax+Bz-c\rVert^2.
\end{aligned}
\]
На допустимом множестве штраф равен нулю, поэтому он не меняет
исходную целевую функцию. При этом квадрат невязки улучшает кривизну
в направлениях, видимых отображением
\((x,z)\mapsto Ax+Bz\), но не создаёт кривизну в его ядре и потому не
делает произвольную невыпуклую задачу глобально сильно выпуклой.

Метод множителей совместно вычисляет
\[
    (x^{k+1},z^{k+1})
    \in
    \operatorname*{argmin}_{x,z}
    \mathcal L_\rho(x,z,y^k)
\]
и затем обновляет
\[
    y^{k+1}
    =
    y^k+\rho(Ax^{k+1}+Bz^{k+1}-c).
\]
Совместная подзадача обычно не разделяется, поскольку раскрытие
квадрата порождает перекрёстный член
\(2x^{\mathsf T}A^{\mathsf T}Bz\). В обычной двойственной
декомпозиции такого штрафа нет, поэтому блоки разделяются, но
теряется регуляризация нарушения ограничения.

\section{Итерация ADMM}

ADMM заменяет совместную минимизацию двумя последовательными шагами:
\[
\begin{aligned}
    x^{k+1}
    &\in
    \operatorname*{argmin}_{x}
    \mathcal L_\rho(x,z^k,y^k),
    \\
    z^{k+1}
    &\in
    \operatorname*{argmin}_{z}
    \mathcal L_\rho(x^{k+1},z,y^k),
    \\
    y^{k+1}
    &=
    y^k+\rho(Ax^{k+1}+Bz^{k+1}-c).
\end{aligned}
\]
После удаления слагаемых, не зависящих от оптимизируемого блока,
прямые подзадачи принимают вид
\[
\begin{aligned}
    x^{k+1}
    \in
    \operatorname*{argmin}_{x}
    \Bigl\{
        f(x)
        +(y^k)^{\mathsf T}Ax
        +\frac{\rho}{2}
        \lVert Ax+Bz^k-c\rVert^2
    \Bigr\},
\end{aligned}
\]
\[
\begin{aligned}
    z^{k+1}
    \in
    \operatorname*{argmin}_{z}
    \Bigl\{
        g(z)
        +(y^k)^{\mathsf T}Bz
        +\frac{\rho}{2}
        \lVert Ax^{k+1}+Bz-c\rVert^2
    \Bigr\}.
\end{aligned}
\]
Каждая задача содержит только один переменный блок, однако второй шаг
использует уже вычисленное значение \(x^{k+1}\). Поэтому классический
двухблочный ADMM является схемой Гаусса--Зейделя, а не параллельной
схемой Якоби. Распределённая реализация возможна за счёт размещения
блоков на разных узлах и внутреннего распараллеливания каждой
подзадачи, но сами два блочных обновления остаются последовательными.

\section{Масштабированная форма}

Введём масштабированную двойственную переменную
\[
    u^k=\frac{y^k}{\rho}.
\]
Для \(r=Ax+Bz-c\) справедливо тождество
\[
    y^{\mathsf T}r+\frac{\rho}{2}\lVert r\rVert^2
    =
    \frac{\rho}{2}\lVert r+u\rVert^2
    -
    \frac{\rho}{2}\lVert u\rVert^2.
\]
Последнее слагаемое не зависит от прямых переменных. Поэтому ADMM
можно записать как
\[
\begin{aligned}
    x^{k+1}
    &\in
    \operatorname*{argmin}_{x}
    \Bigl\{
        f(x)
        +\frac{\rho}{2}
        \lVert Ax+Bz^k-c+u^k\rVert^2
    \Bigr\},
    \\
    z^{k+1}
    &\in
    \operatorname*{argmin}_{z}
    \Bigl\{
        g(z)
        +\frac{\rho}{2}
        \lVert Ax^{k+1}+Bz-c+u^k\rVert^2
    \Bigr\},
    \\
    u^{k+1}
    &=
    u^k+Ax^{k+1}+Bz^{k+1}-c.
\end{aligned}
\]
Если
\[
    r^{k+1}=Ax^{k+1}+Bz^{k+1}-c,
\]
то
\[
    u^{k+1}=u^k+r^{k+1},
    \qquad
    r^{k+1}=u^{k+1}-u^k.
\]
Таким образом, \(u^k\) аккумулирует предыдущие нарушения ограничения,
но обычно сходится не к нулю, а к \(y^\star/\rho\). Поэтому качество
итерации оценивают по невязкам, а не по норме самой масштабированной
переменной.

\section{Условия оптимальности и невязки}

Для выпуклой задачи седловая точка
\((x^\star,z^\star,y^\star)\) удовлетворяет условиям
\[
    0\in\partial f(x^\star)+A^{\mathsf T}y^\star,
\]
\[
    0\in\partial g(z^\star)+B^{\mathsf T}y^\star,
\]
\[
    Ax^\star+Bz^\star-c=0.
\]
Условие оптимальности \(z\)-подзадачи вместе с двойственным
обновлением даёт точное соотношение
\[
    0
    \in
    \partial g(z^{k+1})+B^{\mathsf T}y^{k+1}.
\]
Для \(x\)-подзадачи получаем
\[
\begin{aligned}
    0
    \in{}&
    \partial f(x^{k+1})+A^{\mathsf T}y^{k+1}
    \\
    &-
    \rho A^{\mathsf T}B(z^{k+1}-z^k).
\end{aligned}
\]
Поэтому двойственную невязку определяют как
\[
    s^{k+1}
    =
    \rho A^{\mathsf T}B(z^{k+1}-z^k).
\]
Прямая невязка измеряет нарушение линейного ограничения, а
двойственная --- отклонение от стационарности по блоку \(x\). При
\[
    r^k\to0,
    \qquad
    s^k\to0
\]
условия KKT выполняются асимптотически.

Практический критерий остановки имеет вид
\[
    \lVert r^k\rVert\le\varepsilon_{\mathrm p},
    \qquad
    \lVert s^k\rVert\le\varepsilon_{\mathrm d}.
\]
Если штраф \(\rho\) изменяется, масштабированную переменную следует
пересчитать так, чтобы немасштабированный множитель
\(y=\rho u\) сохранился.

\section{Сходимость и выбор штрафа}

\begin{theorem}[Базовая сходимость ADMM]
Пусть \(f\) и \(g\) --- собственные замкнутые выпуклые функции,
седловая точка обычной функции Лагранжа существует, обе прямые
подзадачи решаются точно и \(\rho>0\) фиксировано. Тогда
\[
    r^k\to0,
\]
\[
    f(x^k)+g(z^k)\to p^\star,
\]
а двойственная последовательность сходится к двойственному решению:
\[
    y^k\to y^\star.
\]
\end{theorem}

Для этих выводов полный ранг \(A\) и \(B\) не требуется. Однако
сходимость самих последовательностей \(x^k\) и \(z^k\),
единственность решений подзадач и единственность прямого предела могут
потребовать сильной выпуклости, коэрцитивности, ранговых условий или
единственности оптимального решения. Классическая теорема не
распространяется автоматически на произвольные невыпуклые задачи и
многоблочные схемы с тремя и более последовательными обновлениями.

Большое \(\rho\) сильнее подавляет прямую невязку, но может ухудшить
обусловленность подзадач. Малое \(\rho\) делает прямые шаги мягче, но
может замедлить достижение допустимости. В отличие от чистого метода
квадратичных штрафов, для выпуклого двухблочного ADMM не требуется
предел \(\rho_k\to\infty\): базовая теория работает при любом
фиксированном \(\rho>0\). Адаптивное изменение штрафа допустимо лишь
в рамках правила, для которого проверены условия сходимости.

Если подзадачи решаются приближённо, ошибки их оптимальности должны
контролироваться. Одно из достаточных условий имеет вид
\[
    \sum_{k=0}^{\infty}\lVert e_x^{k+1}\rVert<\infty,
    \qquad
    \sum_{k=0}^{\infty}\lVert e_z^{k+1}\rVert<\infty,
\]
хотя точные требования зависят от варианта метода.

\section{Разбиение задачи и распределённая реализация}

Преимущество ADMM определяется выбранным разбиением. Оно полезно,
если обе блочные подзадачи существенно проще исходной совместной
задачи. При размещении \(x\) и \(z\) на разных узлах первый узел
вычисляет \(x^{k+1}\), после чего передаёт второму необходимый вклад,
например \(Ax^{k+1}\); второй узел вычисляет \(z^{k+1}\), а затем
обновляется двойственная переменная. Объём сообщений зависит от
структуры \(A\) и \(B\), а последовательность шагов требует
синхронизации.

В частном случае
\[
    A=I,
    \qquad
    B=-I,
    \qquad
    c=0
\]
ограничение имеет вид \(x=z\), а масштабированные обновления
становятся проксимальными:
\[
    x^{k+1}
    =
    \operatorname{prox}_{f/\rho}(z^k-u^k),
\]
\[
    z^{k+1}
    =
    \operatorname{prox}_{g/\rho}(x^{k+1}+u^k).
\]
Поэтому ADMM особенно эффективен, когда проксимальные операторы
отдельных частей вычисляются просто.

По сравнению с двойственным подъёмом ADMM сохраняет квадратичный
штраф, но теряет полную параллельность прямых обновлений. По сравнению
с методом множителей он заменяет одну связанную минимизацию двумя
блочными. Метод не гарантирует хорошую скорость при неудачном
разбиении, не определяет автоматически оптимальное значение \(\rho\)
и требует отдельной теории для невыпуклых, многоблочных и полностью
параллельных модификаций.
